% vim: tw=80 noai
\documentclass[normaltoc,capchap,capsec,times]{abnt}
\usepackage[utf8]{inputenc}
\usepackage[T1]{fontenc}
\usepackage[brazil]{babel}
\usepackage[alf]{abntcite}
\usepackage[ordem=alf]{tabela-simbolos}
\usepackage{url}
\usepackage{graphicx}
\usepackage{listings}
\usepackage{verbatim}
\usepackage{subfigure}
\usepackage{multicol}
\usepackage{framed}
% O comando abaixo define o diretorio onde devem ser colocadas as imagens.
% Neste caso o diretorio é ./imagens/
\graphicspath{ {./imagens/} }
\def\lstlistingname{Listagem}

%%%%%%%%%%%%%%%%%%%%%%%%%%%%%%%%%%%%%%%%%%%%%%%%%%%%
% Dados da monografia
%%%%%%%%%%%%%%%%%%%%%%%%%%%%%%%%%%%%%%%%%%%%%%%%%%%%

\newcommand{\meunome}{José Lucas dos Santos Borges}
\newcommand{\meutitulo}{Construindo um mecanismo para predição de momentos oportunos para interrupção de motoristas utilizando contexto}
\newcommand{\meusubtitulo}{}
\newcommand{\meuano}{2016}
\newcommand{\meuorientador}{Orientador: \profa\ Vaninha Vieira dos Santos}

%%%%%%%%%%%%%%%%%%%%%%%%%%%%%%%%%%%%%%%%%%%%%%%%%%%%

%% O comando \obs aqui definido permite que o autor faca anotacoes na
%% monografia que aparecem no PDF gerado. Para ativar o comando, descomente
%% a primeira linha e comente a segunda.
%% Exemplo de uso: \obs{Preciso melhorar este parágrafo...}

%\newcommand{\obs}[1]{\underline{\textbf{OBSERVAÇÃO}}: \emph{#1}}
\newcommand{\obs}[1]{}

\def\ordfem{\mbox{\raise .35em \hbox{\underline{\scriptsize a}\ }}}
\def\ordmasc{\mbox{\raise .35em \hbox{\underline{\scriptsize o}\ }}}
\def\profa{Prof\ordfem.}

%%%%%%%%%%%%%%%%%%%%%%%%%%%%%%%%%%%%%%%%%%%%%%%%%%%%

\begin{document}

% Capa com Brasão

\begin{titlepage}
   \begin{center}
    	%logotipo
               \includegraphics[scale=0.3]{brasao_ufba} \\
	%\vspace{0.7in}
              \centering{ 
	      \bf{
	      \LARGE{
		\uppercase{UNIVERSIDADE FEDERAL DA BAHIA} \\
 	      }
	      \Large {
                   	\uppercase{INSTITUTO DE MATEMÁTICA} \\
	      }
                   \large {
                       \uppercase{DEPARTAMENTO DE CIÊNCIA DA COMPUTAÇÃO} \\
                  }
              } }
   \end{center}
   \vfill
   \begin{center}
       \bf{
       \large{\uppercase{\meunome}  \\  }
       }
   \end{center}
   \vspace{0.2in}
   \begin{center}
       \bf{
      	 \LARGE{ \uppercase{\meutitulo} } \\
      	 \Large{ \uppercase{\meusubtitulo} }
         \obs{\\ \Large{Esta versão da monografia contém comentários do autor.
          Para removê-los, redefina o comando LaTeX \texttt{obs}.}}
       }
   \end{center}

   \vfill
   \hspace{\stretch{1}}
   \vfill
   \begin{center}
      \normalsize{
          Salvador \\
          \meuano
       }
   \end{center}

\end{titlepage}

%comando abaixo cria uma capa redundante, mas como a capa com brasão foi 
% feita 'manualmente', não faz sentido usar este comando:
%\capa



%\folhaderosto
% o comando acima foi comentado para não criar uma folha de
% rosto redundante, já que ela feita 'manualmente' abaixo

\begin{titlepage}
 \vfill
 \begin{center}
   {\large \uppercase{ \bf{ \meunome\ } } } \\[7cm]
   {\Huge \uppercase{ \bf{ \meutitulo\ } } }\\[1cm]
   \vfill
   \hspace{.45\textwidth} % posicionando a minipage
   \begin{minipage}{.5\textwidth}
     \begin{espacosimples}
       \bf{
	Anteprojeto da Monografia apresentada ao curso de graduação em Sistemas de Informação,
	Departamento de Ciência da Computação, Instituto de Matemática, Universidade Federal da
	Bahia, como requisito parcial para obtenção do grau de Bacharel em Sistemas de Informação. \\
       }
     \end{espacosimples}
     \begin{espacosimples}
       \meuorientador
     \end{espacosimples}
   \end{minipage}
   \vfill
   Salvador \\
   \meuano
 \end{center}
\end{titlepage}


%% As listas a seguir sao opcionais:
%\listadetabelas
\listadesiglas
%\listadesimbolos

\sumario

% O conteudo da monografia esta' nos seguintes arquivos:
\chapter{Objetivo Geral}

O objetivo desta pesquisa é desenvolver um mecanismo que faça predições de momentos oportunos para interrupção de
motoristas, utilizando elementos de contexto.

\chapter{Objetivos Específicos}

\begin{itemize}
  \item Investigar elementos de contexto que indiquem que é possível interromper um motorista em um dado momento
  \item Investigar e aplicar técnicas que podem ser usadas para a predição de interrupções de motoristas
  \item Desenvolver uma solução que implemente os pontos anteriores de forma satisfatória
  \item Implementar a solucão no aplicativo Meu Possante \cite{meupossante}
  \item Avaliar o desempenho do mecanismo de de interrupções
\end{itemize}

\chapter{Justificativa}

Dirigir um veículo pode ser considerado uma tarefa complexa, já que envolve diversas pequenas ações coordenadas e
que dependem tanto de fatores externos quanto internos. Um motorista deve saber se adaptar a diferentes situações,
agir apropriadamente perante cada uma delas e tomar medidas rápidas em caso de emergência.

Neste contexto, a interrupção inoportuna é um perigo que pode ser fatal, pois demanda atenção do motorista em um momento
onde a distração pode levar a um acidente. Um fato que corrobora isto é o de que pessoas levam até 27\% mais tempo para
concluir uma tarefa quando são interrompidas e tendem a cometer o dobro de erros \cite{bailey2006need}. Exemplos de interrupções que podem
ocorrer para motoristas são barulhos repentinos, conversas com passageiros e notificações em seu celular.

Diariamente, pessoas que utilizam qualquer tipo de dispositivo móvel recebem dezenas de notificações \cite{pielot2014situ}. Tais notificações
demandam a atenção do usuário e podem acabar interrompendo uma tarefa que está sendo executada no momento, dado que a
atenção humana é um recurso finito \cite{simon1971designing}.

Interrupções de motoristas já são um grave problema. 10\% dos acidentes com morte nos Estados Unidos em 2014 tem relação
com algum tipo de distração e 13\% deste tipo de distrações estão relacionadas com o uso de celulares e smartphones \cite{distracted2014}.
Diante disto, prevenir este tipo de distração relacionada a dispositivos móveis torna-se uma necessidade.

A utilização de dispositivos móveis enquanto se está dirigindo, apesar de ser proibido no Brasil é liberado em alguns
outros países, como Estados Unidos (em alguns estados como Flórida e Colorado \cite{cellphoneuse}\cite{distracteddriving}) e Suécia \cite{swedendrive}.

Identificando padrões nos dados contextuais que indiquem momentos oportunos para interromper motoristas, é possível
construir uma ferramenta que os notifique apropriadamente a depender do contexto que eles se encontram e que possa ajudar
na diminuição das estatísticas de morte causadas por distração.

\chapter{Revisão da literatura} \label{cap:revisao}

TO DO

\chapter{Metodologia}
\label{ref:metodologia}

\begin{itemize}
  \item Uma revisão da literatura será realizada
  \item Serão identificados momentos oportunos para interrupção de motoristas
  \item Será feita a concepção e definição do escopo do aplicativo
  \item A solução será projetada
  \item A solução será implementada em um aplicativo
  \item A solução será testada e avaliada

\end{itemize}

\chapter{Resultados esperados}

TO DO

\chapter{Cronograma}

\begin{table}[h]
\centering
\resizebox{\textwidth}{!}{%
\begin{tabular}{|r|c|c|c|c|}
\hline
\multicolumn{1}{|c|}{\multirow{2}{*}{Atividade}}                                                       & \multicolumn{2}{c|}{Setembro}                                       & \multicolumn{2}{c|}{Outubro}                                        \\ \cline{2-5}
\multicolumn{1}{|c|}{}                                                                                 & \multicolumn{1}{l|}{1ª quinzena} & \multicolumn{1}{l|}{2ª quinzena} & \multicolumn{1}{l|}{1ª quinzena} & \multicolumn{1}{l|}{2ª quinzena} \\ \hline
Pesquisa de trabalhos relacionados                                                                     & X                                & X                                & X                                &                                  \\ \hline
Revisão da literatura                                                                                  & X                                & X                                & X                                &                                  \\ \hline
\begin{tabular}[c]{@{}r@{}}Revisão da literatura -\\ Conceitos básicos de contexto\end{tabular}        & X                                &                                  &                                  &                                  \\ \hline
Revisão da literatura - Interrupção                                                                    & X                                &                                  &                                  &                                  \\ \hline
Revisão da literatura - Notificação                                                                    &                                  & X                                &                                  &                                  \\ \hline
\begin{tabular}[c]{@{}r@{}}Revisão da literatura - \\ Interrupção de motoristas\end{tabular}           &                                  & X                                &                                  &                                  \\ \hline
\begin{tabular}[c]{@{}r@{}}Revisão da literatura - \\ Sensores relevantes para a proposta\end{tabular} &                                  & X                                &                                  &                                  \\ \hline
Definição do escopo do projeto                                                                         &                                  &                                  & X                                &                                  \\ \hline
Entrega do anteprojeto                                                                                 &                                  &                                  & X                                &                                  \\ \hline
Defesa do PF I                                                                                         &                                  &                                  &                                  & X                                \\ \hline
\end{tabular}%
}
\end{table}


\bibliography{monografia}

\end{document}
